\begingroup
\thispagestyle{plain}

% 국문 논문요약 페이지입니다.
% 이 파일의 역할: 국문 논문요약을 작성합니다.
% 연구 목적, 연구 방법, 주요 결과, 결론을 짧고 압축적으로 요약합니다.
% 논문요약은 논문 전체를 대표하는 글이므로 본문을 모두 쓴 뒤 마지막에 작성하는 것이 좋습니다.
% 제목, 저자명, 학교/전공, 지도교수명처럼 반복 사용되는 정보는
% sub/0-preamble.tex의 \newcommand 변수에서 수정합니다.
% 이 파일에서는 논문요약 본문, 주요어, 하단 제출 문구를 주로 수정합니다.
\begin{center}
\setstretch{1}
\vspace*{10mm}

% 논문요약 제목: 고정 문구입니다.
{\fontsize{14pt}{22.4pt}\selectfont \textbf{논 문 요 약}\par}

\vspace{10mm}
% 논문 제목: sub/0-preamble.tex의 \thesistitleKor 값을 사용합니다.
% 줄간격 133%(표지와 동일 비율 28/21): 16pt x 4/3 = 21.33pt.
{\fontsize{16pt}{21.33pt}\selectfont \thesistitleKor\par}

\vspace{10mm}
% 저자명: sub/0-preamble.tex의 \authorKor 값을 사용하며, \SpacedAuthorName으로 자간을 조정합니다.
{\fontsize{13pt}{20.8pt}\selectfont \SpacedAuthorName{\authorKor}\par}

\vspace{10mm}
% 소속 및 지도교수명: sub/0-preamble.tex의 값을 사용합니다.
{\fontsize{13pt}{20.8pt}\selectfont \universityKor\ \majorKor\par}
{\fontsize{13pt}{20.8pt}\selectfont (지도교수\hspace{2em}\SpacedAuthorName{\advisorKor})\par}
\vspace{15mm}
\end{center}

\begingroup
\setstretch{1.8}
\normalfont\normalsize\selectfont

% 국문 논문요약 본문: 아래 문단들을 논문 내용에 맞게 수정합니다.
\noindent
논문요약은 논문 전체를 완성한 뒤 마지막에 작성하는 것이 좋다.
논문요약에는 연구 목적, 연구 방법, 주요 결과, 결론을 짧고 압축적으로 제시한다.
첫 문장에서는 연구 주제와 목적을 분명히 밝힌다.
이때 연구 대상, 핵심 개념, 수업 또는 평가 맥락, 분석하려는 현상을 함께 제시하면 독자가 논문의 범위를 빠르게 이해할 수 있다.
필요한 경우 실험, 조사, 이론적 분석, 시뮬레이션, 사례 연구, 혼합 연구 등 연구 방법의 성격도 첫 문단 안에서 간단히 밝힌다.

다음으로 연구 방법을 간결하게 요약한다.
연구 참여자 또는 분석 대상, 자료 수집 절차, 사용한 검사 도구나 분석 자료, 분석 방법을 중심으로 쓴다.
선행 연구와 구별되는 연구 설계, 자료 수집 방법, 분석 방법이 있다면 그 특징을 중심으로 제시한다.
다만 방법 자체가 논문의 핵심 기여가 아닌 경우에는 연구 결과를 이해하는 데 필요한 범위에서만 짧게 서술한다.
논문요약에서는 표, 그림, 장 번호를 직접 참조하기보다 본문을 읽지 않아도 이해될 수 있도록 핵심 절차를 문장으로 풀어 쓰는 것이 좋다.

논문요약에서 가장 중요한 부분은 연구 결과와 결론이다.
연구 결과는 추론이나 과도한 해석을 덧붙이기보다 객관적으로 제시한다.
정량 연구에서는 변화의 방향, 집단 간 차이, 통계적으로 확인된 핵심 결과를 중심으로 쓰고, 정성 연구에서는 도출된 범주, 참여자의 반응 양상, 사례에서 확인된 특징을 중심으로 쓴다.
여러 결과를 모두 나열하기보다 연구 문제에 직접 답하는 결과를 우선 제시한다.
결론에서는 연구 목적에 대해 어떤 답을 얻었는지 간략히 밝히고, 그 결과가 과학교육 이론, 수업 설계, 평가, 교사 교육 또는 교육과정 실행에 어떤 의미를 갖는지 한두 문장으로 정리한다.
독자는 논문요약을 통해 연구의 핵심 결과와 의의를 먼저 확인하려 하므로, 논문요약 전체에서 결과와 결론이 충분한 비중을 차지하도록 작성한다.
마지막으로 연구의 제한점이나 후속 연구 방향을 매우 짧게 덧붙일 수 있으나, 논문요약의 대부분을 제한점 설명에 사용하지 않도록 한다.
주요어는 논문 제목에 포함된 핵심 용어, 연구 대상, 연구 방법, 중요한 과학교육 개념을 중심으로 4--6개 정도 제시한다.

\vskip 0.5cm
\noindent
% 주요어: 논문 주제에 맞게 쉼표로 구분하여 수정합니다.
\textbf{주요어:} \thesiskeywordsKor

\vfill

% 하단 제출 문구: 학위수여 연월은 sub/0-preamble.tex의 \ThesisConferralYear, \ThesisConferralMonth에서 수정합니다.
\noindent{\setstretch{1}\fontsize{10pt}{16pt}\selectfont ※ 이 논문은 \ThesisConferralDateKor\ 한국교원대학교 대학원위원회에 제출된 교육학 \degreeProgramKor{}학위 논문임.\par}
\endgroup
\endgroup

\newpage
